\chapter{解析几何}\label{ch:g10:coordgeom}

笛卡尔的伟大思想在于用成对的数描述点，将几何问题
化为计算。仅凭两个公式 ---
中点公式与距离公式 --- 便能证明一个三角形是等腰三角形、
一个四边形是平行四边形，或三点共线，
而无需作出任何一条辅助线。

\section{平面中的坐标}

\begin{definition}[坐标系]\label{def:g10:coordgeom:system}
平面的\emph{坐标系}\index{坐标系}由一个原点 $O$ 与
过 $O$ 的两条带刻度的坐标轴组成：水平的
$x$ 轴与竖直的 $y$ 轴。平面上每一点 $M$ 都有唯一一对
有序数对 $(x, y)$，称为该点的\emph{坐标}\index{坐标}：
即\emph{横坐标}\index{横坐标} $x$ 与
\emph{纵坐标}\index{纵坐标} $y$。
当两轴互相垂直且长度单位相同时，
称该坐标系为\emph{标准正交坐标系}\index{标准正交坐标系}。
本章中的坐标系均为标准正交的。
\end{definition}

\begin{omfigure}
\begin{tikzpicture}[scale=0.85]
  \draw[thin, gray!40] (-2.4,-1.4) grid (4.4,3.4);
  \draw[-Stealth] (-2.6,0) -- (4.6,0) node[below, font=\small] {$x$};
  \draw[-Stealth] (0,-1.6) -- (0,3.6) node[left, font=\small] {$y$};
  \node[below left, font=\small] at (0,0) {$O$};
  \fill[omDef] (3,2) circle (2pt);
  \node[above right, font=\small, omDef] at (3,2) {$M(3, 2)$};
  \draw[dashed, omDef] (3,0) -- (3,2) -- (0,2);
  \fill[omThm] (-2,-1) circle (2pt);
  \node[below left, font=\small, omThm] at (-2,-1) {$N(-2, -1)$};
  \foreach \x in {1,2,3,4} \node[below, font=\footnotesize] at (\x,0) {\x};
  \foreach \y in {1,2,3} \node[left, font=\footnotesize] at (0,\y) {\y};
\end{tikzpicture}

{\small 平面上每一点都由两个数确定：先是
横坐标（水平方向），再是纵坐标（竖直方向）。}
\end{omfigure}

\section{线段的中点}

\begin{proposition}[中点公式]\label{prop:g10:coordgeom:midpoint}
设 $A(x_A, y_A)$ 与 $B(x_B, y_B)$，则
\emph{中点}\index{中点} $I$ 是线段 $[AB]$ 的中点，其坐标为
\[
I\left(\frac{x_A + x_B}{2},\ \frac{y_A + y_B}{2}\right).
\]
\end{proposition}

\begin{proof}
考虑水平方向的坐标。从 $A$ 到 $B$，横坐标
改变量为 $x_B - x_A$；中点的横坐标为
\[
x_A + \frac{x_B - x_A}{2} = \frac{2x_A + x_B - x_A}{2}
= \frac{x_A + x_B}{2},
\]
即两个横坐标的平均数。对纵坐标作同样的计算，
同理得到中点的纵坐标。
\end{proof}

\begin{example}\label{ex:g10:coordgeom:midpoint}
已知 $A(-1, 4)$ 与 $B(5, -2)$，则 $[AB]$ 的中点为
\[
I\left(\frac{-1 + 5}{2},\ \frac{4 + (-2)}{2}\right) = I(2, 1).
\]
公式也可反过来用：若 $A(-1, 4)$ 且其中点为
$I(2, 1)$，则 $B$ 满足 $\frac{-1 + x_B}{2} = 2$ 且
$\frac{4 + y_B}{2} = 1$，于是 $x_B = 5$，$y_B = -2$：点 $B$ 是
$A$ 关于 $I$ 的\emph{中心对称点}。
\end{example}

\section{两点间的距离}

\begin{theorem}[距离公式]\label{thm:g10:coordgeom:distance}
在标准正交坐标系中，$A(x_A, y_A)$ 与
$B(x_B, y_B)$ 之间的距离为
\[
AB = \sqrt{(x_B - x_A)^2 + (y_B - y_A)^2}.
\]
\end{theorem}

\begin{proof}
设点 $C$ 为 $(x_B, y_A)$：它与 $A$ 的纵坐标相同，且与 $B$ 的
横坐标相同，故三角形 $ACB$ 在 $C$ 处为直角。其
两条直角边为水平与竖直方向的线段，长度分别为
$AC = \abs{x_B - x_A}$，$CB = \abs{y_B - y_A}$。由
勾股定理，
\[
AB^2 = AC^2 + CB^2 = (x_B - x_A)^2 + (y_B - y_A)^2,
\]
两边均非负，开平方即得公式。
\end{proof}

\begin{omfigure}
\begin{tikzpicture}[scale=0.85]
  \draw[thin, gray!40] (-0.4,-0.4) grid (5.4,3.4);
  \draw[-Stealth] (-0.6,0) -- (5.6,0) node[below, font=\small] {$x$};
  \draw[-Stealth] (0,-0.6) -- (0,3.6) node[left, font=\small] {$y$};
  \coordinate (A) at (1,1);
  \coordinate (B) at (5,3);
  \coordinate (C) at (5,1);
  \draw[omDef, thick] (A) -- (B);
  \draw[omThm, thick] (A) -- (C) node[midway, below, font=\small]
    {$\abs{x_B - x_A}$};
  \draw[omThm, thick] (C) -- (B) node[midway, right, font=\small]
    {$\abs{y_B - y_A}$};
  \draw[thin] (C) ++(-0.25,0) -- ++(0,0.25) -- ++(0.25,0);
  \fill (A) circle (2pt) node[above left, font=\small] {$A$};
  \fill (B) circle (2pt) node[above, font=\small] {$B$};
  \fill (C) circle (2pt) node[below right, font=\small] {$C$};
\end{tikzpicture}

{\small 距离公式即是将勾股定理应用于
以一条水平直角边与一条竖直直角边为两边的直角三角形 $ACB$。}
\end{omfigure}

\begin{example}\label{ex:g10:coordgeom:distance}
已知 $A(1, 1)$ 与 $B(5, 3)$，则：
\[
AB = \sqrt{(5-1)^2 + (3-1)^2} = \sqrt{16 + 4} = \sqrt{20} = 2\sqrt 5 .
\]
距离通常保留为精确（根式）形式；仅在
最后需要小数答案时，才在最后取近似值。
\end{example}

\begin{remark}
点的顺序无关紧要：$(x_A - x_B)^2 =
(x_B - x_A)^2$。但不要忘记平方！距离
\emph{不是} $\abs{x_B - x_A} + \abs{y_B - y_A}$。
\end{remark}

\section{公式在几何中的应用}

\begin{method}[三角形的形状]\label{met:g10:coordgeom:triangle}
已知三点 $A$、$B$、$C$ 的坐标：
\begin{enumerate}
  \item 利用距离公式计算三个平方长 $AB^2$、$AC^2$、$BC^2$，
        （保留平方可避免开方）；
  \item 若两个平方长相等 $\Rightarrow$ 三角形是
        \emph{等腰三角形}；
  \item 若最大的平方长等于其余两个平方长之和，则
        三角形是\emph{直角三角形}，直角位于最长
        边所对的顶点处，由勾股定理的逆定理得出。
\end{enumerate}
\end{method}

\begin{example}\label{ex:g10:coordgeom:triangle}
设 $A(0, 1)$、$B(4, 3)$ 与 $C(2, -3)$，则
\begin{align*}
AB^2 &= (4-0)^2 + (3-1)^2 = 16 + 4 = 20, \\
AC^2 &= (2-0)^2 + (-3-1)^2 = 4 + 16 = 20, \\
BC^2 &= (2-4)^2 + (-3-3)^2 = 4 + 36 = 40 .
\end{align*}
首先，$AB^2 = AC^2$，故 $AB = AC$，三角形是以 $A$ 为顶点的等腰三角形。
此外 $AB^2 + AC^2 = 40 = BC^2$，由勾股定理的逆定理
可知三角形在 $A$ 处也是直角三角形。
\end{example}

\begin{method}[判定平行四边形]\label{met:g10:coordgeom:parallelogram}
四边形 $ABCD$ 是平行四边形，当且仅当其对角线
$[AC]$ 与 $[BD]$ 有相同的中点。因此：用中点公式算出两条对角线的中点，
并加以比较。
\end{method}

\begin{example}\label{ex:g10:coordgeom:parallelogram}
设 $A(-1, 0)$、$B(2, 2)$、$C(5, 1)$ 与 $D(2, -1)$。
$[AC]$ 的中点为 $\left(\frac{-1+5}{2}, \frac{0+1}{2}\right) = (2, \frac12)$；
$[BD]$ 的中点为 $\left(\frac{2+2}{2}, \frac{2-1}{2}\right) =
(2, \frac12)$，两者重合，故 $ABCD$ 是平行四边形。
\end{example}

\begin{omfigure}
\begin{tikzpicture}[scale=0.75]
  \draw[thin, gray!40] (-1.4,-1.4) grid (5.4,2.4);
  \draw[-Stealth] (-1.6,0) -- (5.6,0);
  \draw[-Stealth] (0,-1.6) -- (0,2.6);
  \coordinate (A) at (-1,0);
  \coordinate (B) at (2,2);
  \coordinate (C) at (5,1);
  \coordinate (D) at (2,-1);
  \draw[omDef, thick] (A) -- (B) -- (C) -- (D) -- cycle;
  \draw[omThm, dashed] (A) -- (C);
  \draw[omThm, dashed] (B) -- (D);
  \fill (A) circle (2pt) node[left, font=\small] {$A$};
  \fill (B) circle (2pt) node[above, font=\small] {$B$};
  \fill (C) circle (2pt) node[right, font=\small] {$C$};
  \fill (D) circle (2pt) node[below, font=\small] {$D$};
  \fill[omThm] (2,0.5) circle (2pt) node[above right, font=\small] {$I$};
\end{tikzpicture}

{\small $ABCD$ 是平行四边形，因为它的两条对角线
有相同的中点 $I(2, \frac12)$。}
\end{omfigure}

\section{练习}

\begin{exercise}[$\star$]\label{exo:g10:coordgeom:1}
设 $A(2, 5)$ 与 $B(-4, 1)$。求 $[AB]$ 的中点
坐标及距离 $AB$。
\end{exercise}

\begin{exercise}[$\star$]\label{exo:g10:coordgeom:2}
设 $A(3, -2)$ 与 $I(1, 2)$。求点 $B$ 的坐标，
使 $I$ 是 $[AB]$ 的中点。
\end{exercise}

\begin{exercise}[$\star$]\label{exo:g10:coordgeom:3}
在点 $P(4, 1)$、$Q(-3, 2)$ 与 $R(0, -5)$ 中，哪一点
离原点 $O(0,0)$ 最近？请以精确距离作答。
\end{exercise}

\begin{exercise}[$\star$]\label{exo:g10:coordgeom:4}
在方格纸上描出点 $A(1, 2)$、$B(5, 4)$、$C(7, 0)$，然后
通过计算说明三角形 $ABC$ 是等腰三角形。它是直角三角形吗？
\end{exercise}

\begin{exercise}[$\star\star$]\label{exo:g10:coordgeom:5}
设 $A(-2, 1)$、$B(1, 3)$、$C(4, 1)$ 与 $D(1, -1)$。
\begin{enumerate}
  \item 证明 $ABCD$ 是平行四边形。
  \item 求 $AB$ 与 $BC$。$ABCD$ 是菱形吗（四边相等）？
  \item 求 $AC$ 与 $BD$。$ABCD$ 是矩形吗（对角线相等）？
\end{enumerate}
\end{exercise}

\begin{exercise}[$\star\star$]\label{exo:g10:coordgeom:6}
圆心为 $\Omega(2, 1)$、半径为 $5$ 的圆由所有
到某定点距离为 $5$ 的点组成，该定点即 $\Omega$。在点 $A(5, 5)$、
$B(-1, -3)$、$C(6, 2)$ 中，哪些在该圆上？哪些在圆内？哪些在圆外？
\end{exercise}

\begin{exercise}[$\star\star$]\label{exo:g10:coordgeom:7}
设 $A(1, 1)$、$B(7, 5)$。求所有点 $M(x, 0)$（即位于 $x$ 轴上的点），
且到 $A$、$B$ 距离相等。（列出 $MA^2 = MB^2$，并解出
$x$。）
\end{exercise}

\begin{exercise}[$\star\star$]\label{exo:g10:coordgeom:8}
设 $A(0, 3)$、$B(4, 1)$。
\begin{enumerate}
  \item 求点 $I$ 的坐标，它是线段 $[AB]$ 的中点。
  \item 证明 $M(2, 2) = I$，再通过计算 $MA$、$MB$
        验证 $M$ 到 $A$、$B$ 距离相等。
  \item 求 $y$ 轴上另一个与 $A$、
        $B$ 距离相等的点。
\end{enumerate}
\end{exercise}

\begin{exercise}[$\star\star$]\label{exo:g10:coordgeom:9}
已知点 $A(-1, -1)$、$B(3, 1)$ 与 $C(11, 5)$。求 $AB$、
$BC$ 与 $AC$，并由此推出 $A$、$B$、$C$ 共线。（三点
共线当且仅当三个距离中最大者等于
另外两者之和。）
\end{exercise}

\begin{exercise}[$\star\star\star$]\label{exo:g10:coordgeom:10}
设 $A(a, 0)$、$B(0, b)$，其中 $a, b > 0$，记 $I$ 为
线段 $[AB]$ 的中点。通过计算证明 $OI = IA = IB$，其中 $O$ 为
原点。（这证明了一个经典定理：直角三角形的
斜边中点到三个顶点等距。）
\end{exercise}

\section{问题：笛卡儿之桥 --- 圆的方程与如何
确定自己的位置}

\begin{problem}[{周末问题 --- 圆化为方程，旧定理
化为计算，三个信标定位一点，
即代数中的三边测量}]
\label{pb:g10:coordgeom:1}
1637 年，笛卡儿在两片大陆之间架起了一座桥梁：
几何中的每一条曲线都变成了 \emph{方程}，每一个
几何问题都化为计算。本题沿着这座桥
双向前行 --- 由距离公式（见\cref{thm:g10:coordgeom:distance}）推导圆的方程，
用三行代数重新证明
经典定理，最后落脚于桥上如今车流最密的一处：
由距离信号计算位置，
也就是 GPS 中“地球视为平面”的核心。

\medskip
\noindent\textbf{第一部分 --- 圆得到方程。}
\begin{enumerate}
  \item 点 $M(x, y)$ 位于圆心为
        $\Omega(2, 1)$、半径为 $5$ 的圆上，当且仅当
        $\Omega M = 5$。将此条件平方，
        得到该圆的方程：
        \[
        (x - 2)^2 + (y - 1)^2 = 25 .
        \]
  \item 检验各点是否属于该圆：在 $A(5, 5)$、$B(6, 4)$、
        $D(4, 4)$ 中，哪些在圆上、哪些在圆内、哪些在圆外？
  \item 一个乔装的圆：用配方法
        （见\cref{pb:g10:algebra:1}）处理下式
        \[
        x^2 + y^2 - 6x + 4y - 12 = 0
        \]
        并给出其圆心与半径。
  \item 对 $x^2 + y^2 - 6x + 4y + 14 = 0$ 做同样的处理。
        配方后的形式揭示了什么？陈述判别准则：何时
        $(x - h)^2 + (y - k)^2 = c$ 表示一个圆、
        一个点，还是什么都不表示？
  \item 将第 3 题中的圆与水平直线
        $y = 3$ 联立：代入并求解。交点有几个，
        这样的直线在几何上叫什么？
\end{enumerate}

\medskip
\noindent\textbf{第二部分 --- 旧定理的三行证明。}
\begin{enumerate}[resume]
  \item 初中分册中的圆周角定理（半圆所对的圆周角是直角），
        用代数重新证明：设 $A(-r, 0)$、$B(r, 0)$ 为
        一条直径的两端，$M(x, y)$ 为圆
        $x^2 + y^2 = r^2$ 上的任意一点。求 $MA^2 + MB^2$，并
        由毕达哥拉斯定理的逆定理得出
        $\widehat{AMB}$ 是直角。
  \item 中线定理：取 $A(-a, 0)$、$B(a, 0)$（于是
        线段 $[AB]$ 的中点即原点 $I$），证明对
        任一点 $M(x, y)$：
        \[
        MA^2 + MB^2 = 2\,MI^2 + \frac{AB^2}{2} .
        \]
  \item 由此求点 $M$ 的轨迹，其中
        $MA^2 + MB^2 = 26$（取 $AB = 6$）：它是什么曲线，
        圆心在哪里，半径为多少？
  \item 另一种轨迹：$A(0,0)$、$B(3,0)$，求所有
        满足 $MA = 2\,MB$。（两边平方、展开、配方：
        一个著名的圆出现了 --- 阿波罗尼奥斯
        不借助坐标就已认识它。）
  \item 用一两句话回答：笛卡儿之桥
        在 \emph{如何} 证明定理上带来了什么改变？将
        第 6 题与初中分册中圆周角定理的矩形证明
        （半圆所对的圆周角是直角）相比较。
\end{enumerate}

\medskip
\noindent\textbf{第三部分 --- 三个信标找到你。}
你的位置 $M(x, y)$ 未知。信标 $P(0, 0)$ 测得
你到它的距离为 $5$；信标 $Q(6, 0)$ 也测得 $5$。
\begin{enumerate}[resume]
  \item 写出两个圆的方程，将它们相减，观察
        平方项相消：剩下的是什么简单方程？
        再与其中一个圆联立，求出 \emph{两个}候选
        位置。
  \item 第三个信标 $R(0, 8)$ 测得你的距离为 $5$。
        计算它到每个候选位置的距离，
        据此判断你所在的位置。
  \item 解释为什么两个圆的方程相减
        \emph{总是}得到一条直线的方程（哪些项
        相消？），以及两圆相交于两点时，这条直线
        在几何上是什么？
  \item 真实的卫星定位发生在空间中，且时钟并不完美：
        每颗卫星（借助信号传播时间）给出一个距离，
        即一个球面。接收机有多少个未知量
        （位置 \emph{和}自身的时钟误差），
        因此 GPS 为何至少要收听
        \emph{四}颗卫星？
  \item 精度问题：设信标 $P$ 的距离实际为
        $5.1$ 而非 $5$（其余不变）。重做第 11 题的
        相减，求出新的 $x$，并
        量化估计值移动了多少。（比较
        \cref{pb:g10:numbers:1}中的误差预算思想。）
\end{enumerate}

\medskip
\noindent\textbf{第四部分 --- 测量员的工具箱。}
\begin{enumerate}[resume]
  \item 最短浇水路径：营地位于 $A(1, 5)$，谷仓
        位于 $B(7, 3)$，一条笔直的河沿坐标轴 $y = 0$ 流过。
        从 $A$ 出发先到河边、再到 $B$，要使步行路程最短：
        作 $B$ 关于河的对称点 $B'$，
        并求线段 $[AB']$ 与河的交点。
        求出交点，并求最短总路程。
        （这一思路即初中分册中的台球反弹练习，
        现在可以完全计算。）
  \item 与另一条路径比较以验证最短性：计算经过
        $Q(3, 0)$ 的总路径，并检验它比你的答案更长。
  \item 判断四边形 $A(1,1)$、$B(4,2)$、
        $C(5,5)$、$D(2,4)$ 属于哪一类四边形：平行四边形？菱形？
        矩形？正方形？
        （参见\cref{met:g10:coordgeom:parallelogram}，
        \cref{met:g10:coordgeom:triangle} --- 也要
        比较对角线。）
  \item 已知 $A(-1, 2)$、$B(3, 4)$、$C(6, 0)$，求点 $D$，使得
        $ABCD$ 为平行四边形（利用对角线中点的
        技巧，即初中分册中关于中心对称的周末问题思路，现用公式表达）。
  \item 总结：三种工具 --- 中点公式、距离公式、
        圆方程相减。对每一种工具，指出它在本题中
        解决了哪一类问题，并
        说明接下来两章会为这套工具增添什么
        （方向与斜率：向量与直线方程）。
\end{enumerate}
\end{problem}
